% !TeX program = lualatex
% !TeX encoding = utf8
% !TeX spellcheck = uk_UA
% !TeX root = ../ElMagConspect.tex

%=========================================================
\chapter[Рух заряджених частинок]{Рух заряджених частинок в електричних і магнітних полях}\label{\currfilebase}
\graphicspath{{\currfilebase/Pictures}}
%=========================================================


%% --------------------------------------------------------
\section{Рух частинок в однорідних полях}
%% --------------------------------------------------------

%% --------------------------------------------------------
\subsection*{Рух в однорідному електричному полі}
%% --------------------------------------------------------

Якщо напруженість поля $\Efield = \const$, то з рівняння руху
\begin{equation}\label{eq:motion_E}
	m\dot{\vect{v}} = q\Efield
\end{equation}
випливає
\begin{equation}\label{eq:motion_E_solution}
	\vect{v} = \vect{v}_0 + \frac{q\Efield}{m}t, \qquad
	\vect{r} = \vect{r}_0 + \vect{v}_0 t + \frac{q\Efield}{2m}t^2,
\end{equation}
тобто має місце рівноприскорений рух з прискоренням, напрямленим уздовж вектора напруженості поля.


%% --------------------------------------------------------
\subsection*{Рух в однорідному магнітному полі}
%% --------------------------------------------------------

На частинку діє сила Лоренца, так що рівняння руху має вигляд
\begin{equation}\label{eq:Lorentz_motion}
	m\dot{\vect{v}} = \frac{q}{c}\vect{v} \times \vect{B}.
\end{equation}
Очевидно, що $\dot{\vect{v}} \perp \vect{v}$, $\dot{\vect{v}} \perp \vect{B}$. Розкладемо вектор швидкості на складові, паралельну і перпендикулярну полю: $\vect{v} = \vect{v}_{\parallel} + \vect{v}_{\perp}$. Для цих складових маємо рівняння
\begin{equation}\label{eq:v_parallel_motion}
	m\dot{\vect{v}}_{\parallel} = 0,
\end{equation}
\begin{equation}\label{eq:v_perp_motion}
	m\dot{\vect{v}}_{\perp} = \frac{q}{c}\vect{v}_{\perp} \times \vect{B}.
\end{equation}
З рівняння \eqref{eq:v_parallel_motion} випливає $\vect{v}_{\parallel} = \const$.

Друге рівняння перепишемо у вигляді
\begin{equation}\label{eq:v_perp_rotation}
	\dot{\vect{v}}_{\perp} = \vect{\omega} \times \vect{v}_{\perp}, \qquad
	\vect{\omega} = -\frac{q\vect{B}}{mc}.
\end{equation}

Воно описує обертання навколо напрямку магнітного поля з кутовою швидкістю $\omega = |q|B/mc$, яка називається \emphx{циклотронною частотою}, --- рис.~\ref{tikz:cyclotron}.

\begin{figure}[h!]
	\centering
	\localinput{CycletronicTrajectory.tikz}
	\caption{Обертальний рух частинки із зарядом $q > 0$ навколо напрямку магнітного поля.}
	\label{tikz:cyclotron}
\end{figure}

Якщо в момент попадання в область дії поля складова швидкості частинки $\vect{v}_{\perp}$ дорівнювала (за величиною) $v_{\perp}$, то радіус кола виявляється рівним
\begin{equation}\label{eq:cyclotron_radius}
	R = \frac{v_{\perp}}{\omega} = \frac{mcv_{\perp}}{|q|B}.
\end{equation}


%% --------------------------------------------------------
\subsection*{Рух в однорідних паралельних електричному та магнітному полях ($\Efield \parallel \vect{B}$)}
%% --------------------------------------------------------


В рівнянні руху
\begin{equation}\label{eq:motion_EB_parallel}
	m\dot{\vect{v}} = q\Efield + \frac{q}{c}\vect{v} \times \vect{B}
\end{equation}
покладемо $\vect{v} = \vect{v}_{\parallel} + \vect{v}_{\perp}$, де компоненти швидкості відповідно $\vect{v}_{\parallel} \parallel \vect{B}$ і $\vect{v}_{\perp} \perp \vect{B}$. Для цих компонент отримуємо рівняння
\begin{equation}\label{eq:v_parallel_E}
	m\dot{\vect{v}}_{\parallel} = q\Efield,
\end{equation}
\begin{equation}\label{eq:v_perp_B}
	m\dot{\vect{v}}_{\perp} = \frac{q}{c}\vect{v}_{\perp} \times \vect{B}.
\end{equation}

\begin{SCfigure}[2][h!]
	\centering
	\localinput{orbit.tikz}
	\caption{Траєкторія частинки в паралельних електричному і магнітному полях.}
	\label{tikz:helix}
\end{SCfigure}

Перше з них описує рівноприскорений рух уздовж напрямку електричного поля, а друге рівняння --- обертальний рух навколо напрямку магнітного поля з циклотронною частотою $\omega = qB/mc$. Траєкторія частинки являє собою спіраль зі зростаючим кроком і з віссю, паралельною силовій лінії магнітного поля (рис.~\ref{tikz:helix}).

%% --------------------------------------------------------
\subsection*{Рух у схрещених полях ($\Efield \perp \vect{B}$)}
%% --------------------------------------------------------

Як і вище, виділимо компоненти швидкості, паралельну та перпендикулярну вектору магнітного поля: $\vect{v} = \vect{v}_{\parallel} + \vect{v}_{\perp}$ (рис.~\ref{tikz:drift}). Для цих компонент маємо рівняння
\begin{equation}\label{eq:v_parallel_crossed}
	m\dot{\vect{v}}_{\parallel} = 0,
\end{equation}
\begin{equation}\label{eq:v_perp_crossed}
	m\dot{\vect{v}}_{\perp} = q\left(\Efield + \frac{1}{c}\vect{v}_{\perp} \times \vect{B}\right).
\end{equation}

Рівняння \eqref{eq:v_parallel_crossed} дає $\vect{v}_{\parallel} = \const$, що означає рух з постійною швидкістю вздовж вектора $\vect{B}$, яка визначається початковими умовами.

\begin{figure}[h!]
	\centering
	\localinput{eldriftplot.tikz}
	\caption{Електричний дрейф додатно (синя крива) та від'ємно (червона крива) заряджених частинок.}
	\label{tikz:drift}
\end{figure}


Для знаходження компоненти $\vect{v}_{\perp}$ покладемо $\vect{v}_{\perp} = \vect{u} + \vect{v}'$ і підберемо вектор $\vect{u}$ так, щоб виключити електричне поле:
\begin{equation}\label{eq:u_condition}
	\Efield + \frac{1}{c}\vect{u} \times \vect{B} = 0.
\end{equation}
У цьому рівнянні можна вважати $\vect{u} \perp \vect{B}$, $\vect{u}\vect{B} = 0$. Для знаходження вектора $\vect{u}$ помножимо почленно рівняння \eqref{eq:u_condition} векторно зліва на $\vect{B}$:
\begin{equation}\label{eq:u_cross_B}
	\vect{B} \times \Efield + \frac{1}{c}\vect{B} \times (\vect{u} \times \vect{B}) = 0.
\end{equation}
Перетворимо подвійний векторний добуток з урахуванням взаємної ортогональності векторів $\vect{B}$ і $\vect{u}$:
\begin{equation}\label{eq:double_cross}
	\vect{B} \times (\vect{u} \times \vect{B}) = \vect{u}B^2 - \vect{B}(\vect{u}\vect{B}) = \vect{u}B^2.
\end{equation}
В результаті отримаємо
\begin{equation}\label{eq:drift_velocity}
	\tcbhighmath{
		\vect{u} = c\frac{\Efield \times \vect{B}}{B^2}.
	}
\end{equation}
Цей рух являє собою дрейф (з постійною швидкістю) в напрямку, перпендикулярному обом полям ($\Efield$ і $\vect{B}$) (рис.~\ref{tikz:drift}). Відповідно для компоненти $\vect{v}'$ отримаємо рівняння
\begin{equation}\label{eq:v_prime_rotation}
	m\dot{\vect{v}}' = \frac{q}{c}\vect{v}' \times \vect{B},
\end{equation}
яке описує обертальний рух з кутовою швидкістю $\omega = -qB/mc$ навколо напрямку магнітного поля.

Сумарний рух являє собою накладення:
\begin{enumerate}
	\item рівномірного руху вздовж вектора $\vect{B}$ зі швидкістю $(\vect{v}_{\parallel})_0$, яка дорівнює компоненті швидкості $\vect{v}_{\parallel}$ в початковий момент,
	\item дрейфового руху зі швидкістю $\vect{u}$,
	\item обертання навколо напрямку $\vect{B}$.
\end{enumerate}

Викладений підхід з виділенням дрейфового руху зі швидкістю $\vect{u}$ справедливий тільки у випадку слабких електричних полів $E \ll B$, коли виконується умова нерелятивістського руху $u = cE/B \ll c$. Якщо ж умова $E \ll B$ не виконується, слід розв'язувати релятивістські рівняння руху.

Знайдемо траєкторію шляхом прямого розв'язання рівнянь \eqref{eq:v_perp_crossed}. Для цього запишемо ці рівняння в системі координат, показаній на рис.~\ref{tikz:drift}.



Введемо орти $\{\vect{i}, \vect{j}, \vect{k}\}$ відповідно вздовж координатних осей $\{x, y, z\}$. Оскільки в розглянутому випадку $\vect{v} \times \vect{B} = \vect{i}B_z v_y - \vect{j}B_z v_x$, то, позначаючи для стислості $E_y = E$, $B_z = B$, маємо:
\begin{equation}\label{eq:system_components}
	\text{(а) } \dot{v}_z = 0; \qquad
	\text{(б) } \dot{v}_x = \frac{qB}{mc}v_y; \qquad
	\text{(в) } \dot{v}_y = \frac{q}{m}E - \frac{qB}{mc}v_x.
\end{equation}

Рівняння \eqref{eq:system_components}а описує рівномірний рух уздовж вектора магнітного поля:
\begin{equation}\label{eq:vz_motion}
	v_z = v_{\parallel} = \const, \qquad z = z_0 + v_{\parallel}t.
\end{equation}
Далі вважаємо $v_{\parallel} = 0$. Для розв'язання рівнянь \eqref{eq:system_components}б і \eqref{eq:system_components}в покладемо
\begin{equation}\label{eq:Vx_substitution}
	V_x = v_x - u, \qquad u = cE/B.
\end{equation}
Після цього рівняння \eqref{eq:system_components}б, в набирають вигляду
\begin{equation}\label{eq:Vx_vy_system}
	\dot{V}_x = \omega v_y, \qquad \dot{v}_y = -\omega V_x,
\end{equation}
де введена циклотронна частота $\omega = qB/mc$ (для $q<0$ вона від'ємна). Виключаючи звідси змінну $V_x$, приведемо цю систему до рівняння гармонічного осцилятора:
\begin{equation}\label{eq:harmonic_oscillator}
	\ddot{v}_y = -\omega^2 v_y,
\end{equation}
звідки знаходимо:
\begin{equation}\label{eq:vy_solution}
	v_y = v_0 \cos(\omega t + \varphi_0), \qquad
	y = y_0 + \frac{v_0}{\omega} \sin(\omega t + \varphi_0).
\end{equation}
За допомогою другого рівняння в \eqref{eq:vy_solution} і \eqref{eq:Vx_substitution} знаходимо $v_x(t)$ і $x(t)$:
\begin{equation}\label{eq:vx_x_solution}
	\begin{gathered}
		v_x = u - \frac{1}{\omega}\dot{v}_y = u + v_0 \sin(\omega t + \varphi_0), \\
		x = x_0 + ut - \frac{v_0}{\omega} \cos(\omega t + \varphi_0).
	\end{gathered}
\end{equation}
Формули \eqref{eq:vy_solution}, \eqref{eq:vx_x_solution} визначають траєкторію частинки за будь-яких початкових умов. Виберемо початок координат так, що $x(0) = 0$, $y(0) = 0$. Тоді
\begin{equation}\label{eq:trochoid}
	\begin{cases}
		x = ut + \frac{v_0}{\omega}[\cos\varphi_0 - \cos(\omega t + \varphi_0)], \\
		y = \frac{v_0}{\omega}[\sin(\omega t + \varphi_0) - \sin\varphi_0],
	\end{cases}
	\qquad
	\begin{cases}
		v_x = u + v_0 \sin(\omega t + \varphi_0), \\
		v_y = v_0 \cos(\omega t + \varphi_0).
	\end{cases}
\end{equation}
Крива, що задається рівняннями \eqref{eq:trochoid}, називається \emphx{трохоїдою}. Її конкретний вигляд залежить від величини $v_0$, яка визначається з початкових умов. Розглянемо частинний випадок, коли початкова швидкість частинки дорівнює нулю:
\begin{equation}\label{eq:zero_initial_velocity}
	v_x(0) = v_y(0) = 0.
\end{equation}
Тоді $\varphi_0 = \pi/2$, $v_0 = -u$, і з \eqref{eq:trochoid} знаходимо:
\begin{equation}\label{eq:cycloid}
	\begin{gathered}
		v_x = u(1 - \cos\omega t), \qquad v_y = u \sin\omega t; \\
		x(t) = ut - \frac{u}{\omega}\sin\omega t, \qquad
		y(t) = \frac{u}{\omega}(1 - \cos\omega t).
	\end{gathered}
\end{equation}

Ця траєкторія являє собою \emphx{циклоїду} (рис.~\ref{tikz:drift_cycloide}). Таку траєкторію описує точка на ободі колеса, що котиться без проковзування по площині. Відповідно до сказаного траєкторія являє собою суперпозицію дрейфу вздовж осі $x$ і обертання навколо осі $z$ з циклотронною частотою $\omega = qB/mc$.

\begin{figure}[h!]
	\centering
	\localinput{eldriftplot_cycloide.tikz}
	\caption{Траєкторія частинки у взаємно перпендикулярних полях $\vect{B}$ і $\Efield$ за умови, що початкова швидкість дорівнює нулю.}
	\label{tikz:drift_cycloide}
\end{figure}

%% --------------------------------------------------------
\subsection*{Ефект Холла}
%% --------------------------------------------------------

Класичний приклад руху заряджених частинок у схрещених полях реалізується у твердому провіднику чи напівпровіднику, по якому тече електричний струм і який знаходиться в магнітному полі, --- це \emphx{ефект Холла}, відкритий Е.~Холлом у 1879~р.

Розглянемо провідну пластинку завтовшки $d$ і завширшки $w$, вздовж якої (вісь $x$) тече струм густиною $\vect{j}$, а перпендикулярно площині пластинки (вісь $z$) прикладене магнітне поле $\vect{B}$ (рис.~\ref{tikz:hall_effect}). Носії заряду --- концентрацією $n$ і зарядом $q$ --- рухаються вздовж провідника з дрейфовою швидкістю $\vect{v}_d$, такою що
\begin{equation}\label{eq:hall_current_density}
	\vect{j} = nq\vect{v}_d.
\end{equation}

\begin{figure}[h!]
	\centering
	\localinput{halleffect.tikz}
%    \includegraphics[width=0.5\linewidth]{HallEffect}
	\caption{Ефект Холла: струм $\vect{j}$ уздовж пластинки, магнітне поле $\vect{B}$ перпендикулярно їй; сила Лоренца відхиляє носії заряду до одного з країв, у результаті чого встановлюється поперечне холлівське поле $\Efield_H$.}
	\label{tikz:hall_effect}
\end{figure}

На кожен носій заряду, що рухається з дрейфовою швидкістю $\vect{v}_d$ в магнітному полі, діє сила Лоренца $\frac{q}{c}\vect{v}_d \times \vect{B}$, яка відхиляє носії до одного з бічних країв пластинки. Внаслідок цього на протилежних гранях накопичується нескомпенсований заряд, що створює поперечне (холлівське) електричне поле $\Efield_H$. У стаціонарному режимі результуюча поперечна сила, що діє на носії, дорівнює нулю, тобто виконується та сама умова \eqref{eq:u_condition}, яку ми використали для знаходження дрейфової швидкості:
\begin{equation}\label{eq:hall_balance}
	q\Efield_H + \frac{q}{c}\vect{v}_d \times \vect{B} = 0
	\qquad \Longrightarrow \qquad
	\Efield_H = -\frac{1}{c}\vect{v}_d \times \vect{B}.
\end{equation}
Отже, холлівське поле є прямим наслідком того самого механізму компенсації електричної та магнітної сил, який лежить в основі дрейфового руху \eqref{eq:drift_velocity}, --- тільки тут роль \enquote{зовнішнього} електричного поля відіграє поле $\Efield_H$, що виникає самоузгоджено.

Виразивши дрейфову швидкість через густину струму за формулою \eqref{eq:hall_current_density}, $v_d = j/(nq)$, знаходимо величину холлівського поля:
\begin{equation}\label{eq:hall_field}
	E_H = \frac{1}{nqc}\,jB.
\end{equation}
Різниця потенціалів між бічними гранями пластинки (холлівська напруга) дорівнює $U_H = E_H h$, а оскільки повний струм $I = jhd$, отримуємо
\begin{equation}\label{eq:hall_voltage}
	\tcbhighmath{
		U_H = \frac{1}{nqc}\,\frac{IB}{d}.
	}
\end{equation}
Коефіцієнт
\begin{equation}\label{eq:hall_coefficient}
	R_H = \frac{1}{nqc}
\end{equation}
називається \emphx{сталою Холла}. Ця формула справедлива для одного типу носіїв заряду з однаковим часом релаксації; для напівпровідників з електронами й дірками вираз для $R_H$ складніший.

Знак холлівської напруги визначається знаком носіїв заряду: він протилежний для провідників з електронною ($q<0$) і дірковою ($q>0$) провідністю, що дозволяє експериментально встановити тип основних носіїв у зразку. За відомими $I$, $B$, $d$ і виміряним $U_H$ формула \eqref{eq:hall_voltage} дає змогу визначити концентрацію носіїв заряду $n$, а вимірювання ефекту Холла широко використовується для характеризації напівпровідників, а також у датчиках магнітного поля (датчиках Холла).


%% --------------------------------------------------------
\subsection*{Дрейф у плазмі та плазмові двигуни}
%% --------------------------------------------------------

Розглянутий електричний дрейф має фундаментальне значення для фізики плазми: оскільки швидкість дрейфу не залежить від знаку заряду й маси частинки, замагнічені електрони та іони в схрещених полях дрейфують \emphx{в один і той самий бік} з однаковою швидкістю. Це принципово відрізняє $\vect{E}\times\vect{B}$-дрейф від руху частинки під дією самого лише електричного поля, який розділяє заряди різного знаку і породжує струм.

Схрещені поля лежать в основі роботи \emphx{холлівських двигунів} (\emphx{Hall-effect thrusters}), одного з найпоширеніших типів електричних ракетних двигунів для космічних апаратів. Проте тут істотною є не однаковість дрейфу електронів та іонів, а, навпаки, різниця в ступені їхнього замагнічення: електрони замагнічені, а іони --- ні.

\paragraph{Принцип дії холлівського двигуна.}
У кільцевому каналі двигуна створюється радіальне магнітне поле $\vect{B}$ та осьове електричне поле $\Efield$, причому $\Efield \perp \vect{B}$. Електрони, інжектовані з катода, потрапляючи в область сильного радіального поля, замагнічуються (їхній ларморівський радіус $r_L = mcv_\perp/(|q|B)$ значно менший за розміри каналу) і, згідно з формулою \eqref{eq:drift_velocity}, здійснюють азимутальний дрейф --- це і є \enquote{\emphx{холлівський струм}}, що дав назву двигуну.

\begin{SCfigure}[1.2][h!]
	\centering
%	\localinput{halleffectthruster.tikz}
    \includegraphics[height=8cm]{Hall_Truster}
	\caption{Схема холлівського двигуна.}
	\label{tikz:hall_thruster}
\end{SCfigure}

Іони ж, маючи значно більший ларморівський радіус (через велику масу), практично не замагнічуються в тому самому полі й вільно прискорюються електричним полем уздовж осі каналу, створюючи тягу. Радіальне магнітне поле, замагнічуючи електрони, різко обмежує їхню рухливість упоперек силових ліній, тобто вздовж осі до анода, тому основне падіння потенціалу зосереджується в зоні прискорення поблизу виходу каналу. Електрони, що дрейфують, одночасно іонізують ударами робочу речовину (зазвичай ксенон), а плазма в каналі в цілому залишається квазінейтральною.


%\paragraph{Магнітне утримання та дрейфи в термоядерній плазмі.}
%Аналогічний $\vect{E}\times\vect{B}$-дрейф відіграє ключову роль в установках магнітного утримання плазми (токамаки, стеларатори). Тут, окрім електричного дрейфу, суттєвими є й інші дрейфи того самого типу --- \emphx{градієнтний дрейф} (при неоднорідності $\vect{B}$) і \emphx{дрейф кривини} магнітних силових ліній, які отримуються аналогічним чином, якщо в рівнянні \eqref{eq:u_condition} замінити електричну силу $q\Efield$ на відповідну ефективну силу:
%\begin{equation}\label{eq:general_drift}
%	\vect{u}_{\vect{F}} = \frac{c}{q}\,\frac{\vect{F} \times \vect{B}}{B^2}.
%\end{equation}
%Зокрема, для градієнтного дрейфу $\vect{F} = -\mu\nabla B$ (де $\mu = mv_\perp^2/2B$ --- магнітний момент частинки), а для дрейфу кривини $\vect{F} = mv_\parallel^2 \vect{R}_c/R_c^2$, де $\vect{R}_c$ --- радіус-вектор кривини силової лінії, напрямлений від центра кривини до точки, де розглядається частинка. Оскільки ці дрейфи, на відміну від електричного, \emphx{залежать від знаку заряду}, вони спричиняють поляризацію плазми та є однією з причин турбулентного переносу, що обмежує якість утримання в термоядерних реакторах.
%
%Таким чином, простий приклад руху зарядженої частинки в схрещених полях виявляється не лише навчальною задачею класичної електродинаміки, а й основою для розуміння як інженерних застосувань (електричні ракетні двигуни), так і фундаментальних проблем керованого термоядерного синтезу.


%% --------------------------------------------------------
\subsection*{Загальний випадок схрещених полів}
%% --------------------------------------------------------

У загальному випадку є складова вектора електричного поля, паралельна магнітному полю: $\Efield_{\parallel} \neq 0$. Тоді для компоненти швидкості $\vect{v}_{\parallel}$ має місце рівняння
\begin{equation}\label{eq:v_parallel_general}
	m\dot{\vect{v}}_{\parallel} = q\Efield_{\parallel}.
\end{equation}
Це означає, що на рух, описаний у попередньому пункті, накладається прискорений рух у напрямку $\Efield_{\parallel} \parallel \vect{B}$. Дрейф зі швидкістю $\vect{u} = c\Efield \times \vect{B}/B^2$ зберігається, так само як і обертання з циклотронною частотою $\omega = -qB/mc$.


%% --------------------------------------------------------
\section{Методи вимірювання питомого заряду електрона}
%% --------------------------------------------------------

\emphx{Питомим зарядом} називається відношення $q/m$, де $m$ --- маса частинки.

Рух частинки за наявності тільки електричних і магнітних полів описується рівнянням
\begin{equation}\label{eq:q_m_motion}
	m\frac{d\vect{v}}{dt} = q\left(\Efield + \frac{1}{c}\vect{v} \times \vect{B}\right).
\end{equation}
Звідси видно, що воно повністю визначається питомим зарядом $q/m$, а не зарядом і масою окремо.


%% --------------------------------------------------------
\subsection*{Вимірювання $q/m$ за відхиленням траєкторії в електричному полі}
%% --------------------------------------------------------

У цьому методі пучок частинок пропускається між обкладками зарядженого конденсатора, і за виміряним відхиленням траєкторії визначають відношення $q/m$. Схема досліду показана на рис.~\ref{tikz:condenser_deflection}.

\begin{figure}[h!]
	\centering
	\localinput{In_condensator.tikz}
	\caption{Відхилення зарядженої частинки електричним полем конденсатора.}
	\label{tikz:condenser_deflection}
\end{figure}

В конденсаторі є однорідне електричне поле з напруженістю $\Efield$, напрямлене перпендикулярно початковій швидкості частинки $\vect{v}_0$. Це поле викликає рух з прискоренням
\begin{equation}\label{eq:a_E}
	\vect{a} = \frac{q}{m}\Efield, \qquad \vect{a} \perp \vect{v}_0.
\end{equation}
Якби поле було відсутнє, то частинки попадали б у точку «0» екрана Е (рис.~\ref{tikz:condenser_deflection}). Під дією поля відхилення частинки в конденсаторі у вертикальному напрямку становить
\begin{equation}\label{eq:y1_E}
	y_1 = \frac{at^2}{2} = \frac{qE}{2m}t^2 = \frac{q}{2m}E\frac{l^2}{v_0^2},
\end{equation}
де $t = l/v_0$ --- час прольоту конденсатора. Тут враховано, що оскільки поле $\Efield$ напрямлене вздовж осі $y$, то компонента швидкості $v_x = \const = v_0$.

На виході з конденсатора швидкість $v_y$ становитиме
\begin{equation}\label{eq:vy_E}
	v_y = at = \frac{qE}{m}\frac{l}{v_0}.
\end{equation}
Далі частинка летить прямолінійно, причому траєкторія утворює кут $\alpha$ з віссю $x$ такий, що
\begin{equation}\label{eq:tg_alpha}
	\tg\alpha = \frac{v_y}{v_x} = \frac{qE}{m}\frac{l}{v_0^2}; \qquad
	y_1 = \frac{l}{2}\tg\alpha.
\end{equation}

Долітаючи до екрана, вона набуває додаткового відхилення
\begin{equation}\label{eq:y2_E}
	y_2 = L \tg\alpha.
\end{equation}
Повне зміщення частинки в результаті становитиме
\begin{equation}\label{eq:y_total_E}
	\tcbhighmath{
		y = y_1 + y_2 = \left(\frac{1}{2}l + L\right)\tg\alpha = \frac{qE}{mv_0^{2}}l\left(\frac{1}{2}l + L\right).
	}
\end{equation}
Таким чином, за виміряним відхиленням частинки $y$ за допомогою формули \eqref{eq:y_total_E} визначається питомий заряд частинки $q/m$.


%% --------------------------------------------------------
\subsection*{Вимірювання $q/m$ за відхиленням траєкторії в магнітному полі}
%% --------------------------------------------------------

Для вимірювання питомого заряду можна замість електричного поля використати магнітне, також напрямлене перпендикулярно траєкторії частинки (рис.~\ref{tikz:magnetic_deflection}). На частинку діє сила Лоренца $\vect{F} = \frac{q}{c}\vect{v} \times \vect{B}$, створюючи прискорення
\begin{equation}\label{eq:a_B}
	\vect{a} = \frac{\vect{F}}{m} = \frac{q}{mc}\vect{v} \times \vect{B}.
\end{equation}

\begin{figure}[h!]
	\centering
	\localinput{In_magfield.tikz}
	\caption{Відхилення зарядженої частинки магнітним полем.}
	\label{tikz:magnetic_deflection}
\end{figure}

Нехай швидкість частинки велика, і відхилення за час прольоту поля мале. Тоді можна прийняти
\begin{equation}\label{eq:a_B_approx}
	v_x \approx v_0, \qquad
	\vect{a} \approx \frac{q}{mc}\vect{v}_0 \times \vect{B}, \qquad
	a = |\vect{a}| \approx \frac{q}{mc}v_0 B.
\end{equation}
Відхилення частинки за час прольоту поля $t = l/v_0$ становитиме
\begin{equation}\label{eq:y1_B}
	y_1 = \frac{1}{2}at^2 = \frac{q}{2m}B\frac{l^2}{v_0 c}.
\end{equation}
Швидкість $v_y$ частинки на виході з області дії поля рівна
\begin{equation}\label{eq:vy_B}
	v_y = at = \frac{ql}{mc}B,
\end{equation}
а кут, який утворює вектор швидкості з віссю $x$, визначається рівністю
\begin{equation}\label{eq:tg_beta}
	\tg\beta = \frac{v_y}{v_x} \approx \frac{v_y}{v_0} = \frac{qlB}{mv_0 c}, \qquad
	y_1 = \frac{l}{2}\tg\beta.
\end{equation}
Долітаючи до екрана, частинка отримує додаткове зміщення
\begin{equation}\label{eq:y2_B}
	y_2 = L \tg\beta.
\end{equation}
Повне зміщення частинки становитиме
\begin{equation}\label{eq:y_total_B}
	\tcbhighmath{
		y = y_1 + y_2 = \left(\frac{1}{2}l + L\right)\tg\beta = \frac{qB}{mv_0 c}l\left(\frac{1}{2}l + L\right).
	}
\end{equation}
Вимірюючи це зміщення, можна знайти питомий заряд $q/m$ частинки.


%% --------------------------------------------------------
\subsection*{Метод Томсона}
%% --------------------------------------------------------

Дж.~Дж.~Томсон використав компенсаційну схему, в якій рух частинки управлявся як електричним, так і магнітним полем. Ця схема дозволяла не вимірювати попередньо швидкість частинок $v_0$. Суть методу в наступному. Спочатку Томсон включав тільки магнітне поле (як описано вище) і вимірював відхилення пучка заряджених частинок:
\begin{equation}\label{eq:Thomson_B}
	y = \frac{qB}{mv_0 c}l\left(\frac{1}{2}l + L\right).
\end{equation}
Потім він додатково включав електричне поле, перпендикулярне магнітному ($\Efield \perp \vect{B}$), і підбирав його величину таким чином, щоб повністю компенсувати відхилення пучка, повернув його в точку $y = 0$. Умовою компенсації була рівність
\begin{equation}\label{eq:Thomson_compensation}
	qE = \frac{q}{c}v_0 B, \qquad \text{або} \qquad v_0 = c\frac{E}{B}.
\end{equation}
Тому наведений вище вираз для зміщення пучка в магнітному полі переписується у вигляді
\begin{equation}\label{eq:Thomson_y}
	y = \frac{q}{m}\frac{B^2}{c^2 E}l\left(\frac{1}{2}l + L\right).
\end{equation}
Звідси за відомими значеннями зміщення $y$ і полів $B$ та $E$ визначається питомий заряд $q/m$.

Описаним методом Дж.~Дж.~Томсон в 1897~р. при вивченні катодних променів вперше виміряв питомий заряд електрона. Сучасне значення питомого заряду електрона (Томсон отримав величину того ж порядку) становить
\begin{equation}\label{eq:e_m_value}
	e/m = -5{,}273 \cdot 10^{17} \ \text{од.~СГСЕ} \cdot \text{г}^{-1}
	= -1{,}759 \cdot 10^{11} \ \text{Кл} \cdot \text{кг}^{-1}.
\end{equation}


%% --------------------------------------------------------
\section{Досліди Міллікена}
%% --------------------------------------------------------

У попередніх дослідах вимірювався питомий заряд частинок. Р.~Міллікен в 1909~р. запропонував спосіб безпосереднього вимірювання заряду. Схема його установки показана на рис.~\ref{tikz:Millikan}.

У простір між пластинами конденсатора Міллікен вводив за допомогою пульверизатора найдрібніші крапельки масла. Краплі іонізувалися, і їхнім рухом можна було керувати, змінюючи величину і знак заряду на пластинах конденсатора.

\begin{figure}[h!]
	\centering
	\includegraphics[width=0.75\linewidth]{Milliken}
	\caption{Схема установки Міллікена для вимірювання заряду електрона (сила $q\Efield$ вказана для випадку $q < 0$).}
	\label{tikz:Millikan}
\end{figure}

Розглянемо метод докладніше. Будемо для визначеності вважати, що заряд краплі від'ємний, $q < 0$. У формулах нижче, де йдеться про величини сил і швидкостей, під $q$ розуміємо $|q|$.

На краплю діють такі сили:
\begin{enumerate}
	\item сила з боку електричного поля $q\Efield$;
	\item рівнодійна сили тяжіння і архімедової сили
	\begin{equation}\label{eq:gravity_force}
		\vect{F}_{\text{т}} = \frac{4\pi}{3}r^3(\rho - \rho_0)\vect{g},
	\end{equation}
	де $r$ --- радіус краплі, $\rho$ і $\rho_0$ --- густини краплі і повітря відповідно;
	\item сила в'язкого тертя
	\begin{equation}\label{eq:viscous_force}
		\vect{F}_{\text{гр}} = -6\pi\eta r\vect{v}.
	\end{equation}
\end{enumerate}

За відсутності електричного поля крапля падає з постійною швидкістю $v_0$, яка визначається з умови $F_{\text{т}} = F_{\text{гр}}$, або
\begin{equation}\label{eq:falling_equilibrium}
	\frac{4\pi}{3}r^3(\rho - \rho_0)g = 6\pi\eta r v_0.
\end{equation}

Увімкнувши електричне поле, змушують краплю підніматися з невеликою швидкістю $v_E$, яку можна знайти з рівняння
\begin{equation}\label{eq:rising_equilibrium}
	\frac{4\pi}{3}r^3(\rho - \rho_0)g + 6\pi\eta r v_E = qE.
\end{equation}
З урахуванням рівності \eqref{eq:falling_equilibrium} перепишемо це співвідношення у вигляді
\begin{equation}\label{eq:charge_balance}
	6\pi\eta r(v_0 + v_E) = qE,
\end{equation}
звідки можна знайти заряд краплі:
\begin{equation}\label{eq:charge_drop}
	q = 6\pi\eta r\frac{v_0 + v_E}{E}.
\end{equation}

Практично вимірювання здійснювалося таким способом. За відсутності електричного поля вимірювалася швидкість вільного падіння краплі $v_0$. Потім із співвідношення \eqref{eq:falling_equilibrium} визначався радіус краплі:
\begin{equation}\label{eq:drop_radius}
	r = \sqrt{\frac{9\eta v_0}{2g(\rho - \rho_0)}}.
\end{equation}
Увімкнувши електричне поле і вимірявши швидкість підйому краплі $v_E$, за формулами \eqref{eq:charge_drop}--\eqref{eq:drop_radius} або
\begin{equation}\label{eq:charge_full}
	q = 6\pi\eta\sqrt{\frac{9\eta v_0}{2g(\rho - \rho_0)}}\frac{v_0 + v_E}{E}
	= 9\pi\sqrt{\frac{2\eta^3 v_0}{g(\rho - \rho_0)}}\frac{v_0 + v_E}{E}
\end{equation}
можна знайти заряд $q$. Слід мати на увазі, що для крапель мікронних розмірів закон Стокса потребує поправки (Каннінгема) на проковзування газу, оскільки радіус краплі порівнянний з довжиною вільного пробігу молекул повітря; Міллікен враховував таку поправку в обробленні результатів.

\begin{Attention}
В дослідах Міллікена мікроскопічні краплі олії набували електричного заряду внаслідок двох послідовних процесів: початкової трибоелектризації через тертя об сопло розпилювача та подальшого захоплення іонів з газового середовища. У робочій камері приладу повітря піддавали штучній іонізації рентгенівським випромінюванням, яке вибивало електрони з нейтральних молекул газу, утворюючи вільні електрони та позитивно заряджені аероіони. Під час руху в хмарі іонізованого газу дисперсні краплі захоплювали ці носії заряду, причому дискретний характер адсорбції зумовлював стрибкоподібну зміну сумарного заряду краплі на величину, чітко кратну елементарному електричному заряду, що й дозволило експериментально довести дискретність електрики та обчислити заряд електрона.
\end{Attention}


Аналізуючи результати експериментів, Міллікен встановив, що зміни заряду крапель $\Delta q$ і сам заряд $q$ були завжди цілими кратними одній і тій же мінімальній величині:
\begin{equation}\label{eq:elementary_charge}
	e \approx 4{,}80 \cdot 10^{-10} \ \text{Фр} = 1{,}602 \cdot 10^{-19} \ \text{Кл}.
\end{equation}
За відомими значеннями питомого заряду $e/m$ і заряду $e$ електрона можна знайти масу:
\begin{equation}\label{eq:electron_mass}
	\frac{e}{e/m} = m \approx 9{,}11 \cdot 10^{-28} \ \text{г}.
\end{equation}
Значення досліду Міллікена полягає в тому, що була експериментально доведена дискретність електричного заряду.

\begin{Questions}
\begin{quizlist}
	\item Запишіть рівняння руху зарядженої частинки в однорідному електричному полі та його розв'язок. Яким буде рух і куди напрямлене прискорення?
	\item Чому сила Лоренца не виконує роботи і не змінює модуль швидкості частинки? Як розкладають швидкість на складові $\vect{v}_{\parallel}$ і $\vect{v}_{\perp}$ в магнітному полі та який рух відповідає кожній із них?
	\item Що таке циклотронна частота? Від яких величин вона залежить, а від яких (наприклад, від швидкості частинки) --- ні? Запишіть формулу для радіуса кола $R$.
	\item Якою є траєкторія частинки в однорідних паралельних полях $\Efield \parallel \vect{B}$ (рис.~\ref{tikz:helix})? Чому крок спіралі зростає?
	\item У схрещених полях $\Efield \perp \vect{B}$ введено швидкість $\vect{u}$ (рис.~\ref{tikz:drift}). З якої умови її знаходять, і чому їй відповідає формула $\vect{u} = c\,\Efield\times\vect{B}/B^2$? Куди напрямлений дрейф відносно $\Efield$ і $\vect{B}$?
	\item Чому швидкість електричного дрейфу не залежить від знака заряду й маси частинки? Чому при цьому вона все ж не поділяє заряди різного знака, на відміну від руху під дією самого лише електричного поля?
	\item З яких трьох рухів складається рух у схрещених полях? За якої умови справедливий такий підхід і що робити, якщо вона не виконується?
	\item Що таке трохоїда та циклоїда? За якої початкової умови траєкторія в схрещених полях є циклоїдою (рис.~\ref{tikz:drift_cycloide}) і з чим її можна порівняти механічно?
	\item У чому полягає ефект Холла (рис.~\ref{tikz:hall_effect})? Чому в стаціонарному режимі результуюча поперечна сила на носії дорівнює нулю?
	\item Виведіть формулу для холлівської напруги $U_H$ і сталої Холла $R_H$. Які величини можна визначити за виміряною $U_H$ і як за знаком холлівської напруги визначити тип носіїв?
	\item Поясніть принцип дії холлівського двигуна. Чому електрони в каналі замагнічені, а іони --- ні, і як це забезпечує тягу?
	\item Що таке питомий заряд? Чому рух частинки в полях визначається саме відношенням $q/m$, а не зарядом і масою окремо?
	\item Опишіть метод вимірювання $q/m$ за відхиленням в електричному полі конденсатора (рис.~\ref{tikz:condenser_deflection}). З яких двох частин складається повне зміщення $y = y_1 + y_2$?
	\item У чому перевага методу Томсона порівняно з вимірюванням окремо в електричному та магнітному полі? Яка умова компенсації відхилення і що з неї випливає для швидкості $v_0$?
	\item Опишіть установку Міллікена (рис.~\ref{tikz:Millikan}). Які сили діють на краплю і як за швидкостями $v_0$ та $v_E$ визначають її радіус і заряд?
	\item Що довели досліди Міллікена і чому саме зміни заряду крапель були кратні одній і тій же величині? За відомими $e$ та $e/m$ як знайти масу електрона?
\end{quizlist}
\end{Questions}